\newif\ifglobal

%%%%%%%%%%%%%%%%%%%%%%%
%%% Compile to view %%%
%%%%%%%%%%%%%%%%%%%%%%%

%%%%%%%%%%%%%%%%%%%%%%%%%%%%%%%%%%%%%%%%%%%%%%%%%%%
%%% Readme:                                     %%%
%%% https://www.overleaf.com/read/bwdjvfjksvjy/ %%%
%%%%%%%%%%%%%%%%%%%%%%%%%%%%%%%%%%%%%%%%%%%%%%%%%%%

% >>> M?: marks a module setting number or important notice

% >>> M4: Set {./relative-path} to external files for this module <<<
\def \modulefiles {./03-SYSMEM}
% To add an external file, use:
% (Replace '---' with actual filename)
% '\input{\modulefiles/tables/---.tex}' for tables
% '\includegraphics{\modulefiles/figures/---}' for figures
% '\subfile{\modulefiles/---}' for register files

% >>> M1: This module is in global project? <<<
% 'YES' - uncomment; 'NO' - comment out
\globaltrue

\ifglobal
    \documentclass[main\_\_EN.tex]{subfiles}
   
\else
    \documentclass[openany, 10pt]{book}

    % >>> M2: In ./00-shared/preamble-trm-module, also find
    %         settings for visibility of todo notes, watermark <<<
    \usepackage{preamble-trm-module}

    % >>> M3: Temporary labels in English,
    %         you can add your own labels there <<<
    \myexternaldocument{./00-shared/temp-labels-en}

\fi %\ifglobal


\setcounter{secnumdepth}{4}


% >>> M5: If this module needs any updates in
% './00-shared/preamble-trm-module.sty'
% (include additional package, etc.), contact your module PM
% or create the file './00-shared/changelog.tex' make the
% necessary updates yourself and document them in this file <<<

% Make text left-aligned and allow latex to break pages naturally, comment out for full justification
\raggedright
\raggedbottom


\begin{document}


% Sets language into which pre-defined bits of text are translated
\selectlanguage{english}
\setenglishformatting

% Do not delete this block
\ifglobal
% Add the following front matter if
% this module is outside of global project
\else
    \InputIfFileExists{readme}{}{}
    \listoftodos
    {\let\clearpage\relax \listofdonetodos}
    \tableofcontents
    %\listoftables
    %\listoffigures
\fi



%%%%%%%%%%%%%%%%%%%%%%%%%
%%% TRM Module Begins %%%
%%%%%%%%%%%%%%%%%%%%%%%%%

\hypertarget{sysmem}{}
\chapter{System and Memory}
\label{mod:sysmem}


\section{Overview}

The ESP32-S2 is a single-core system with one Harvard Architecture Xtensa$^®$ LX7 CPU. All internal memory, external memory, and peripherals are located on the CPU buses.

\section{Features}

\begin{itemize}
    \item {\bfseries Address Space}
        \begin{itemize}
            \item 4 GB (32 bits wide) address space in total accessed from the data bus and instruction bus
    	    \item 464 KB internal memory address space accessed from the instruction bus
    	    \item 400 KB internal memory address space accessed from the data bus
    	    \item 1.77 MB peripheral address space
    	    \item 7.5 MB external memory virtual address space accessed from the instruction bus
    	    \item 14.5 MB external memory virtual address space accessed from the data bus
    	    \item 320 KB internal DMA address space
    	    \item 10.5 MB external DMA address space
        \end{itemize}
    \item {\bfseries Internal Memory}
        \begin{itemize}
            \item 128 KB Internal ROM
	        \item 320 KB Internal SRAM
	        \item 8 KB RTC FAST Memory
	        \item 8 KB RTC SLOW Memory
	    \end{itemize}
    \item {\bfseries External Memory}
        \begin{itemize}
            \item Supports up to 1 GB external SPI flash
	        \item Supports up to 1 GB external SPI RAM
	    \end{itemize}
    \item {\bfseries DMA}
        \begin{itemize}
            \item 9 DMA-supported modules / peripherals
        \end{itemize}
\end{itemize}

Figure \ref{figure:System_AddressMapping_Structure} illustrates the system structure and address mapping.

\begin{figure}[H]
    \centering
    \includegraphics[width=1.0\textwidth]{\modulefiles/figures/newchip_memorymapping_20191118}
    \caption{System Structure and Address Mapping}
    \label{figure:System_AddressMapping_Structure}
\end{figure}

\begin{tiplisting}
\vspace{-2em}
\begin{itemize}
\item The memory space with gray background is not available to users.
\item The range of addresses available in the address space may be larger or smaller than the actual available memory of a particular type.
\end{itemize}
\end{tiplisting}

\section{Functional Description}

\subsection{Address Mapping}

The Harvard Architecture Xtensa$^®$ LX7 CPU can address 4 GB (32 bits wide) memory space.

Addresses below \hexprefix{}4000\_0000 are serviced using the data bus. Addresses in the range \hexprefix{}4000\_0000 \verb+~+ \hexprefix{}4FFF\_FFFF are serviced using the instruction bus. Addresses over and including \hexprefix{}5000\_0000 are shared by both data and instruction bus.

Both data bus and instruction bus are little-endian. The CPU can access data via the data bus in a byte-, half-word-, or word-aligned manner. The CPU can also access data via the instruction bus, but only in a word-aligned manner; non-word-aligned access will cause a CPU exception.

The CPU can:
\begin{itemize}
    \item directly access the internal memory via both data bus and instruction bus;
    \item access the external memory which is mapped into the address space via cache;
    \item access modules / peripherals via data bus.
\end{itemize}

Table \ref{table:Address_Mapping} lists the address ranges on the data bus and instruction bus and their corresponding target memory.

Some internal and external memory can be accessed via both data bus and instruction bus. In such cases, the same memory is available to the CPU at two address ranges.

\subfile{\modulefiles/Address_Mapping_en.tex}



\subsection{Internal Memory}\label{para:internal memory}

The internal memory consists of four segments: Internal ROM (128 KB), Internal SRAM (320 KB), RTC FAST Memory (8 KB), and RTC SLOW Memory (8 KB).

The Internal ROM is broken down into two parts: Internal ROM 0 (64 KB) and Internal ROM 1 (64 KB).

The Internal SRAM is broken down into two parts: Internal SRAM 0 (32 KB) and Internal SRAM 1 (288 KB).

RTC FAST Memory and RTC SLOW Memory are both implemented as SRAM.

Table \ref{table:Embedded_Memory_Address_Mapping} lists all types of internal memory and their address ranges on the data bus and instruction bus.

\subfile{\modulefiles/Embedded_Memory_Address_Mapping_en.tex}



\subsubsection{Internal ROM 0}

Internal ROM 0 is a 64-KB, read-only memory space, addressed by the CPU on the instruction bus via range(s) described in Table \ref{table:Embedded_Memory_Address_Mapping}.


\subsubsection{Internal ROM 1}

Internal ROM 1 is a 64-KB, read-only memory space, addressed by the CPU on the data or instruction bus via range(s) described in Table \ref{table:Embedded_Memory_Address_Mapping}.

The two address ranges access Internal ROM 1 in the same order, so, for example, addresses \hexprefix{}3FFA\_0000 and \hexprefix{}4001\_0000 access the same word, \hexprefix{}3FFA\_0004 and \hexprefix{}4001\_0004 access the same word, \hexprefix{}3FFA\_0008 and \hexprefix{}4001\_0008 access the same word, etc.


\subsubsection{Internal SRAM 0}\label{Internal SRAM 0}

Internal SRAM 0 is a 32-KB, read-and-write memory space, addressed by the CPU on the data or instruction bus, in the same order, via range(s) described in Table \ref{table:Embedded_Memory_Address_Mapping}.

Hardware can be configured to use 8 KB, 16 KB, 24 KB, or the entire 32 KB space in this memory to cache external memory. The space used as cache cannot be accessed by the CPU, while the remaining space can still be accessed by the CPU.


\subsubsection{Internal SRAM 1}

Internal SRAM 1 is a 288-KB, read-and-write memory space, addressed by the CPU on the data or instruction bus, in the same order, via range(s) described in Table \ref{table:Embedded_Memory_Address_Mapping}.

\hypertarget{Trace}{}
Internal SRAM 1 comprises eighteen 16-KB (sub)memory blocks. One block can be used as Trace Memory, in which case this block's address range cannot be accessed by the CPU.

\subsubsection{RTC FAST Memory}
\hypertarget{PRO_CPU_Only}{}
RTC FAST Memory is an 8-KB, read-and-write SRAM, addressed by the CPU on the data or instruction bus, in the same order, via range(s) described in Table \ref{table:Embedded_Memory_Address_Mapping}.

\subsubsection{RTC SLOW Memory}

RTC SLOW Memory is an 8-KB, read-and-write SRAM, addressed by the CPU via range(s) shared by the data bus and the instruction bus, as described in Table \ref{table:Embedded_Memory_Address_Mapping}.

RTC SLOW Memory can also be used as a peripheral addressable to the CPU via either \hexprefix{}3F42\_1000 \verb+~+ \hexprefix{}3F42\_2FFF or \hexprefix{}6002\_1000 \verb+~+ \hexprefix{}6002\_2FFF on the data bus.


\subsection{External Memory} \label{para:external_Memory}

ESP32-S2 supports multiple QSPI/OSPI flash and RAM chips. It also supports hardware encryption/decryption based on XTS-AES to protect user programs and data in the flash and external RAM.

\subsubsection{External Memory Address Mapping}

The CPU accesses the external flash and RAM via the cache. According to the MMU settings, the cache maps the CPU’s address to the external physical memory address. Due to this address mapping, the ESP32-S2 can address up to 1 GB external flash and 1 GB external RAM.

Using the cache, ESP32-S2 can support the following address space mappings at the same time.
  \begin{itemize}
  \item  Up to 7.5 MB instruction bus address space can be mapped into the external flash or RAM as individual 64 KB blocks, via the instruction cache (ICache). Byte (8-bit), half-word (16-bit) and word (32-bit) reads are supported.

  \item  Up to 4 MB read-only data bus address space can be mapped into the external flash or RAM as individual 64 KB blocks, via ICache. Byte (8-bit), half-word (16-bit) and word (32-bit) reads are supported.

  \item Up to 10.5 MB data bus address space can be mapped into the external RAM as individual 64 KB blocks, via DCache. Byte (8-bit), half-word (16-bit) or word (32-bit) reads and writes are supported. Blocks from this 10.5 MB space can also be mapped into the external flash or RAM, for read operations only.
\end{itemize}

Table \ref{table:External_Memory_Address_Mapping} lists the mapping between the cache and the corresponding address ranges on the data bus and instruction bus.

\subfile{\modulefiles/External_Memory_Address_Mapping_en.tex}


\subsubsection{Cache}

As shown in Figure \ref{figure:Cache_arch}, the caches on ESP32-S2 are separated and allow prompt response upon simultaneous requests from the data bus and instruction bus. Some internal memory space can be used as cache (see Section \ref{Internal SRAM 0}). When a cache miss occurs, the cache controller will initiate a request to the external memory. When ICache and DCache simultaneously initiate a request, the arbiter determines which gets the access to the external memory first. The cache size of ICache and DCache can be configured to 8 KB and 16 KB, respectively, while their block size can be configured to 16 bytes and 32 bytes, respectively.

\begin{figure}[H]
    \centering
    \fbox{\includegraphics[width=0.45\textwidth]{\modulefiles/figures/Cache_arch_en}}
     \caption{Cache Structure}
    \label{figure:Cache_arch}
\end{figure}


\subsubsection{Cache Operations}
ESP32-S2 caches support the following operations:
\begin{enumerate}
\item \textbf{Invalidate}: The cache clears the valid bit of a tag. The CPU needs to access the external memory in order to read/write the data. There are two types of invalidation operations: manual invalidation and automatic invalidation. Manual invalidation performs only on data in the specified area in the cache, while automatic invalidation performs on all data in the cache. Both ICache and DCache have this function.

\item \textbf{Clean}: The cache clears the dirty bit of the tag and retains the valid bit. The CPU can then read/write the data directly from the cache. Only DCache has this function.

\item \textbf{Write-back}: The cache clears the dirty block flag of the tag and retains the valid bit. It also forces the data in the corresponding address to be written back to the external memory. The CPU can then read/write the data directly from the cache. Only DCache has this function.

\item \textbf{Preload}: To preload a cache is to load instructions and data into the cache in advance. The minimum unit of a preloading is one block. There are two types of preloading: manual preloading and automatic preloading. Manual preloading means that the hardware prefetches a piece of continuous data according to the virtual address specified by the software. Automatic preloading means the hardware prefetches a piece of continuous data according to the current hit / miss address (depending on configuration).

\item \textbf{Lock / Unlock}: There are two types of lock: prelock and manual lock. When prelock is enabled, the cache locks the data in the specified area when filling the missing data to cache memory. When manual lock is enabled, the cache checks the data that has been filled into the cache memory and locks the data that falls in the specified area. The data in the locked area is always stored in the cache and will not be replaced. But when all the ways within the cache are locked, the cache will replace data, as if the ways were not locked. Unlocking is the reverse of locking. The manual invalidation, clean, and write-back operations are only available after unlocking.

\end{enumerate}


\subsection{DMA}
With DMA, the ESP32-S2 can perform data transfers between:
\begin{itemize}
\item modules / peripherals and internal memory;
\item different types of internal memory;
\item modules / peripherals and external memory;
\item internal and external memory.
\end{itemize}

DMA uses the same addressing as the data bus to read and write Internal SRAM 0 and Internal SRAM 1. Specifically, DMA uses address range \hexprefix{}3FFB\_0000 \verb+~+ \hexprefix{}3FFB\_7FFF to access Internal SRAM 0, and \hexprefix{}3FFB\_8000 \verb+~+ \hexprefix{}3FFF\_FFFF to access Internal SRAM 1. Note that DMA cannot access the internal memory occupied by the cache.

In addition, DMA addresses the external RAM from \hexprefix{}3F50\_0000 \verb+~+ \hexprefix{}3FF7\_FFFF, the same used by the CPU to access DCache. When DCache and DMA access the external memory simultaneously, data consistency is required.

Nine modules / peripherals on the ESP32-S2 support DMA, as shown in Table \ref{table:Module_with_DMA}. With DMA, some of them can only access internal memory, some can access both internal and external memory.

For more information on DMA, please refer to Chapter \ref{mod:dma}: \textit{DMA Controller}.

\subfile{\modulefiles/Module_with_DMA_en.tex}

\subsection{Modules/Peripherals}

The CPU can access modules/peripherals via address range \hexprefix{}3F40\_0000 \verb+~+ \hexprefix{}3F4F\_FFFF on the data bus, or via \hexprefix{}6000\_0000 \verb+~+ \hexprefix{}600B\_FFFF and \hexprefix{}6180\_0000 \verb+~+ \hexprefix{}6180\_3FFF shared by the data bus and instruction bus.

\subsubsection{Naming Conventions for Peripheral Buses} \label{Convention: Peripheral Bus Name}
There are two peripheral buses defined as follows:
\begin{itemize}
    \item \textbf{PeriBus1}\label{PeripheralBus1}: Which refers to the address range \hexprefix{}3F40\_0000 \verb+~+ \hexprefix{}3F4F\_FFFF on the bus. \hexprefix{}3F40\_0000 is the base address.
    \item \textbf{PeriBus2}\label{PeripheralBus2}: Which refers to the address ranges \hexprefix{}6000\_0000 \verb+~+ \hexprefix{}600B\_FFFF and \hexprefix{}6180\_0000 \verb+~+ \hexprefix{}6180\_3FFF on the bus. \hexprefix{}6000\_0000 is the base address.
\end{itemize}

All references to ``\hyperref[PeripheralBus1]{PeriBus1}" and ``\hyperref[PeripheralBus2]{PeriBus2}" in this document indicate the corresponding address range(s).

\subsubsection{Differences Between \hyperref[PeripheralBus1]{PeriBus1} and \hyperref[PeripheralBus2]{PeriBus2}} \label{differences_between_buses}

The CPU can access modules / peripherals more efficiently through \hyperref[PeripheralBus1]{PeriBus1} than through \hyperref[PeripheralBus2]{PeriBus2}. However, \hyperref[PeripheralBus1]{PeriBus1} features speculative reads, which means it cannot guarantee that each read is valid. Therefore, the CPU has to use \hyperref[PeripheralBus2]{PeriBus2} to access some special registers, for example, FIFO registers.

In addition, \hyperref[PeripheralBus1]{PeriBus1} will upset the order of r/w operations on the bus to improve performance, which may cause programs that have strict requirements on the r/w order to crash. In such cases, please add volatile before the program statement, or use \hyperref[PeripheralBus2]{PeriBus2} instead.

\subsubsection{Module / Peripheral Address Mapping}

Table \ref{table:Peripheral_Address_Mapping} lists all the modules / peripherals and their respective address ranges. Note that addresses in column ``Boundary Address" are offsets relative to the base address, instead of absolute addresses. The absolute addresses are the addition of bus base address and the corresponding offsets.

\subfile{\modulefiles/Peripheral_Address_Mapping_en.tex}


\subsubsection{Addresses with Restricted Access from \hyperref[PeripheralBus1]{PeriBus1}} \label{special register inside peripherals}

As mentioned in Section \ref{differences_between_buses}, \hyperref[PeripheralBus1]{PeriBus1} features speculative reads, which means it is forbidden to read FIFO registers. Table \ref{table:Special_Registers_for_Dport_Bus} below lists the address (range) with restricted access from \hyperref[PeripheralBus1]{PeriBus1}.

There are four reserved user-defined registers that can be configured as needed to add more addresses with restricted access. For more information, please refer to Chapter \ref{mod:permctrl}: \textit{Permission Control}.

\subfile{\modulefiles/Special_Registers_for_Dport_Bus_en.tex}

\end{document}
